\documentclass[12pt]{article}
\usepackage[utf8]{inputenc}
\usepackage[T1]{fontenc}
\usepackage{hyperref}
\title{Amatl\'an Zapotec Language Resources for K-12 Curriculum Development}
\author{Compiled by Claude Sonnet 5.5 (claude-sonnet-5-5) for Andr\'as Kornai}
\date{}
\begin{document}
\maketitle

\section{Introduction}
Amatl\'an Zapotec (ISO 639-3 zpo, Glottocode amat1238) is an Otomanguean language spoken in southern Oaxaca, Mexico, in San Crist\'obal Amatl\'an and San Francisco Logueche, with about 10,000 speakers. Search found mainly Bible translation and audio materials plus reference pages; no school materials were found.

\section{Primary Educational Resources}
No primers, textbooks or school programmes were found in search results.

\section{Linguistic and Reference Materials}
\begin{itemize}
\item Wikipedia article on Amatl\'an Zapotec (it cites an archived SIL Mexico page by Sylvia Riggs, whose content I could not verify). \url{https://wikipedia.com/wiki/ISO_639:zpo}
\item Glottolog entry (amat1238), linking to Endangered Languages Project and OLAC records. \url{https://glottolog.org/resource/languoid/id/amat1238}
\item Wiktionary category for Amatl\'an Zapotec. \url{https://en.wiktionary.org/wiki/Category:Amatl%C3%A1n_Zapotec_language}
\end{itemize}

\section{Religious and Community Resources}
\begin{itemize}
\item Scripture Earth: New Testament (PDF), read-along and audio viewer, a Spanish Old Testament selection, and a link to an Android Bible app with text and audio. \url{https://www.scriptureearth.org/00eng.php?iso=zpo}
\item Global Recordings Network: evangelistic and Bible-teaching audio, and a link to a Faith Comes By Hearing New Testament recording. \url{https://globalrecordings.net/en/language/zpo}
\end{itemize}

\section{Digital and Audio Resources}
The Scripture Earth and Global Recordings Network items above include audio (downloadable zip files and a slideshow package). All of it is religious in content.

\section{Curriculum Development Context}
The language is used as a first language by the whole community but is not known to be taught in schools. The New Testament text and audio are the largest bodies of written and recorded Amatl\'an Zapotec found, and they could serve as a corpus for orthography and vocabulary work. Secular teaching materials would have to be created with community speakers.

\end{document}